Dukes' 1D classification establishes the framework on which the present
2D analysis rests. In that work, the coin operator is parameterised by
a real amplitude $\rho$ and a relative phase $\delta = \alpha + \beta$,
and the condition for a $k$-cycle to admit a full-state revival
$U_k^N = I_{2k}$ is that every eigenphase be a rational multiple of
$2\pi$. Dukes enumerated solutions for $k \in \{2,3,4,6\}$ up to
$N = 30$ and identified a finite set of solutions for $k \in
\{5,8,10\}$, but did not prove completeness of these lists; no exact
solutions are known for other cycle lengths. The revival set
$\{2,3,4,5,6,8,10\}$ that governs our 2D analysis is therefore an
established result, whereas the specific 1D period sets remain only
partially characterised. Section~\ref{sec:results} presents the 2D
enumeration directly, without relying on an assumed 1D period list.
